\renewcommand{\Anno}{2007}

\renewcommand{\NumeroProblema}{3}

\renewcommand{\Testo}{
La Scuola S.\ Anna organizza da anni l'operazione ``Turisti d'Eccellenza'': chiunque, pagando 300 euro, può trascorrere 24 ore insieme a docenti e studenti della Scuola.  Per il momento, al massimo una persona al giorno è ammessa a questa attività: per questo motivo occorre prenotare per tempo e pagare in anticipo. Inspiegabilmente, negli ultimi anni è accaduto spesso che alcuni aspiranti turisti, pur avendo regolarmente prenotato e pagato, non si sono poi presentati alla giornata designata, anche se questo ha comportato per loro la perdita della quota di partecipazione.
	
Elaborando i dati acquisiti, gli statistici hanno calcolato che il turista che ha pagato e prenotato si presenta con una certa probabilità $p$.

Per minimizzare i posti vuoti e massimizzare i profitti, si è deciso da quest'anno di seguire una politica di \emph{overbooking}. La Scuola pertanto, pur avendo un solo posto disponibile per ogni giorno, si riserva di venderne anche 2 o 3 (non di più, perché una legge regionale lo vieta).  Come sempre, chi non si presenta perde la quota di partecipazione già versata; se invece in un dato giorno arrivano in 2 o più, solo uno trascorrerà la giornata alla Scuola (chi arriva primo), mentre gli altri avranno diritto alla restituzione della quota versata più un indennizzo pari a 500 euro.
	
Determinare, in funzione di $p$, la strategia ottimale per la Scuola, e cioè il numero di posti da mettere in vendita per ogni giorno (1, 2 o 3) in modo da assicurarsi, in media, il maggior guadagno possibile.
}

\renewcommand{\Risposta}{
La situazione in funzione di $p$ è la seguente.
\begin{itemize}
	\item Se $0\leq p<1/2$, allora alla Scuola conviene mettere in vendita 3 posti al giorno.

	\item Se $1/2<p<\frac{1}{2}\sqrt{3/2}$, allora alla Scuola conviene mettere in vendita 2 posti al giorno.

	\item Se $\frac{1}{2}\sqrt{3/2}<p\leq 1$, allora alla Scuola conviene mettere in vendita un solo posto al giorno.
	
	\item Nel caso limite $p=1/2$ è indifferente mettere in vendita 2 o 3 posti. Nel caso limite $p=\frac{1}{2}\sqrt{3/2}$ è indifferente mettere in vendita 1 o 2 posti.
\end{itemize}
}

\renewcommand{\Aiutino}{
Indichiamo con $G_{k}(p)$ il guadagno medio della Scuola quando questa mette in vendita $k$ posti (con $k=1,2,3$). Dimostrare che
$$G_{1}(p)=3,
\hspace{3em}
G_{2}(p)=6-8p^{2},
\hspace{3em}
G_{3}(p)=9-24p^{2}+8p^{3}.$$

Resta da calcolare, al variare di $p$, il massimo tra le 3 funzioni. Per far questo occorre confrontare $G_{1}$ e $G_{2}$, quindi $G_{2}$ e $G_{3}$. Il confronto tra $G_{1}$ e $G_{3}$ non serve.
}

\renewcommand{\SoluzioneUfficiale}{Calcoliamo il guadagno medio per la Scuola nel caso in cui metta in
vendita 1, 2 o 3 posti per giorno. Per comodità esprimiamo i prezzi in centinaia di euro.

Se la Scuola mette in vendita un solo posto, allora non vi è incertezza, ed il guadagno medio è
$$G_{1}(p)=3.$$

Se la Scuola mette in vendita 2 posti, allora incassa 6, a cui bisogna togliere
\begin{itemize}
	\item 8 (3 di restituzione e 5 di indennizzo) nel caso in cui i turisti si presentino entrambi, cosa che accade con probabilità $p^{2}$;

	\item  0 nel restante caso, cioè se i turisti presentano in 0 o 1.
\end{itemize}

Ne segue che il guadagno medio è dato da
$$G_{2}(p)=6-8p^{2}.$$

Se la Scuola mette in vendita 3 posti, allora incassa 9, a cui bisogna togliere
\begin{itemize}
	\item  16 (cioè $8\times 2$) se si presentano in 3, cosa che accade con probabilità $p^{3}$;

	\item 8 se si presentano in 2, cosa che accade con probabilità $3p^{2}(1-p)$ (infatti $p^{2}$ è la probabilità che in 2 si presentino, $1-p$ la probabilità che il restante non si presenti, 3 sono i modi di scegliere quello che non si presenta);
	
	\item  0 nei restanti casi.
\end{itemize}

Ne segue che in questo caso il guadagno medio è dato da
$$G_{3}(p)=9-16p^{3}-8\cdot 3p^{2}(1-p)=9-24p^{2}+8p^{3}.$$

Non resta quindi che stabilire, al variare di $p$, quale delle 3
funzioni è più grande. Ricordando che comunque $0\leq p\leq
1$, abbiamo che
$$G_{2}(p)\geq G_{1}(p)\ \Longleftrightarrow\ 8p^{2}\leq 3\ \Longleftrightarrow\ 
p\leq \frac{1}{2}\sqrt{\frac{3}{2}},$$
con uguaglianza se e solo se $p=\frac{1}{2}\sqrt{3/2}$. Inoltre
$$G_{3}(p)\geq G_{2}(p)\ \Longleftrightarrow\ 8p^{3}-16p^{2}+3\geq 0
\ \Longleftrightarrow\ (2p-1)(4p^{2}-6p-3)\geq 0.$$

Non è difficile vedere che $4p^{2}-6p-3<0$ per ogni $p\in[0,1]$ (trattandosi di una parabola rivolta verso l'alto basta in realtà controllare la disuguaglianza per $p=0$ e $p=1$) e quindi tutto dipende dal fattore di primo grado $2p-1$.  Ne segue che
$$G_{3}(p)\geq G_{2}(p)\ \Longleftrightarrow\ p\leq\frac{1}{2},$$
con uguaglianza se e solo se $p=1/2$.

Siamo dunque pronti per trarre le conclusioni.
\begin{itemize}
	\item  Se $0\leq p<1/2$, allora $G_{3}(p)>G_{2}(p)>G_{1}(p)$. Ne segue che alla Scuola conviene mettere in vendita 3 posti al giorno.

	\item  Se $1/2<p<\frac{1}{2}\sqrt{3/2}$, allora $G_{2}(p)>G_{3}(p)$ e $G_{2}(p)>G_{1}(p)$. Ne segue che alla Scuola conviene mettere in vendita 2 posti al giorno.

	\item Se $\frac{1}{2}\sqrt{3/2}<p\leq 1$, allora $G_{1}(p)>G_{2}(p)>G_{3}(p)$.  Ne segue che conviene mettere in vendita un solo posto al giorno.
	
	\item Nel caso limite $p=1/2$ è indifferente mettere in vendita 2 o 3 posti. Nel caso limite $p=\frac{1}{2}\sqrt{3/2}$ è indifferente mettere in vendita 1 o 2 posti.
\end{itemize}

Il seguente grafico rappresenta le funzioni $G_{1}(p)$, $G_{2}(p)$, $G_{3}(p)$ al variare di $p$ in $[0,1]$.  La linea spessa rappresenta il massimo tra le 3 funzioni, cioè il guadagno medio della Scuola scegliendo la strategia ottimale in funzione di $p$.

\begin{center}
	\psset{xunit=30ex,yunit=2ex}
	\pspicture(-0.2,-7)(2.2,11)
	\psplot[linecolor=blue]{0}{1}{3}
	\psplot[linecolor=red]{0}{1}{6 8 x 2 exp mul sub}
	\psplot[linecolor=green]{0}{1}{9 24 x 2 exp mul neg 8 x 3 exp mul add add}
	\psline[linewidth=0.5\pslinewidth]{->}(-0.2,0)(1.2,0)
	\psline[linewidth=0.5\pslinewidth]{->}(0,-7)(0,11)
	\psplot[linewidth=2\pslinewidth]{0}{0.5}{9 24 x 2 exp mul neg 8 x 3 exp mul add add}
	\psplot[linewidth=2\pslinewidth]{0.5}{0.5 1.5 sqrt mul}{6 8 x 2 exp mul sub}
	\psplot[linewidth=2\pslinewidth]{0.5 1.5 sqrt mul}{1}{3}
	\psdots(0.5,0)(!0.5 1.5 sqrt mul 0)(1,0)
	\psline[linecolor=magenta,linestyle=dashed](0.5,0)(0.5,4)
	\psline[linecolor=magenta,linestyle=dashed](!0.5 1.5 sqrt mul 0)(!0.5 1.5 sqrt mul 
	3)
	\psline[linecolor=magenta,linestyle=dashed](1,0)(1,3)
	\rput(-0.05,3){3}
	\rput(-0.05,6){6}
	\rput(-0.05,9){9}
	\rput[tl](1,-0.5){1}
	\rput[t](1.2,-0.5){$p$}
	\rput[t](0.5,-0.5){$\displaystyle{\frac{1}{2}}$}
	\rput[t](!0.5 1.5 sqrt mul -0.5){$p_{*}$}
	\psline[linecolor=green](1.4,9)(1.6,9)
	\psline[linecolor=red](1.4,6)(1.6,6)
	\psline[linecolor=blue](1.4,3)(1.6,3)
	\psline[linewidth=2\pslinewidth](1.4,0)(1.6,0)
	\rput[l](1.65,9){$G_{3}(p)$}
	\rput[l](1.65,6){$G_{2}(p)$}
	\rput[l](1.65,3){$G_{1}(p)$}
	\rput[l](1.65,0){$\max\{G_{1}(p),G_{2}(p),G_{3}(p)\}$}
	\rput[l](1.65,-4){$p_{*}=\displaystyle{\frac{1}{2}\sqrt{\frac{3}{2}}}$}
	\endpspicture
\end{center}

Come era intuitivo prevedere, più la probabilità $p$ è bassa, più è facile che i turisti non si presentino, dunque più è conveniente aumentare il numero di posti venduti.
}

\renewcommand{\SoluzioniAlternative}{
}

\renewcommand{\Osservazioni}{\TitoloOsservazione 
Nella soluzione presentata non abbiamo avuto bisogno di confrontare $G_{3}(p)$ e $G_{1}(p)$. Tale confronto produce una disequazione di terzo grado senza radici razionali, dunque non risolubile per via del tutto elementare.

Sarebbe molto interessante generalizzare il problema ammettendo che la Scuola possa mettere in vendita $k$ posti (o ancora più in generale che per ogni giorno ci siano $m$ posti disponibili e $k$ messi in vendita). Intuitivamente si direbbe che più $p$ è piccola e più è conveniente aumentare il numero di posti in vendita. Tuttavia una dimostrazione di questi fatti sembra tutt'altro che ovvia. Il problema è che le (dis)equazioni coinvolte hanno un grado sempre maggiore all'aumentare di $k$ (e di $m$).
}

\renewcommand{\Commenti}{
}